%
% 16-rational.tex -- 
%
% (c) 2026 Prof Dr Andreas Müller
%
\subsection{Rationale Funktionen
\label{buch:analytisch:potenzreihen:subsection:rational}}
Eine rationale Funktion ist ein Quotient der Form
\begin{equation}
f(z)
=
\frac{p(z)}{q(z)},
\qquad
p(z),q(z)\in\mathbb{C}[z]
\label{buch:analytisch:potenzreihen:bruch}
\end{equation}
von zwei Polynomen in $z$.
Polynome sind holomorphe Funktionen, die in ganz $\mathbb{C}$
definiert sind.
Nach der Quotientenregel ist auch der Quotient $f(z)$ eine holomorphe
Funktion überall dort, wo $q(z)$ nicht verschwindet.
In diesem Abschnitt untersuchen wir, wie sich für eine solche Funktion
eine Potenzreihenentwicklung finden lässt.

%
% Polynomieller und rationaler Teil
%
\subsubsection{Polynomieller und rationaler Teil}
Durch Polynomdivision lässt sich der Bruch
\eqref{buch:analytisch:potenzreihen:bruch}
in ein Polynom $s(z)$ und einen einfacheren Bruch
\[
f(z) = s(z) + \frac{r(z)}{q(z)}
\]
zerlegen.
Der Grad von $r(z)$ ist kleiner als der Grad von $q(z)$.
Das Polynom $s(z)$ ist bereits eine Potenzreihe mit unendlichem
Konvergenzradius, es muss nur noch der rationale Teil $r(z)/q(z)$
als Potenzreihe dargestellt werden.

%
% Partialbrüche
%
\subsubsection{Partialbrüche}
In der Analysis lernt man, dass eine rationale Funktion eine
Partialbruchzerlegung hat.
Sie wird dort für die Integration rationaler Funktionen verwendet und
wird daher im Abschnitt~\ref{buch:ratint:bernoulli:subsection:partialbruch}
im Kapitel~\ref{chapter:ratint} über das Integrationsproblem genauer
aufgearbeitet.
Für eine Faktorisierung
\[
q(z) = (z-\alpha_1)^{n_1}\cdot\ldots\cdot (z-\alpha_k)^{n_k}
\]
in Linearfaktoren kann der Bruch als 
\[
\frac{r(z)}{q(z)}
=
\sum_{k=1}^n\sum_{l=1}^{n_k} \frac{A_k^l}{(z-\alpha_k)^l}
\]
geschrieben werden.
Man muss also nur noch eine Potenzreihendarstellung für jeden
einzelnen Summanden finden. 

Ein Summand mit Nennergrad 1 hat die geometrische Reihe
\begin{align*}
\frac{A}{z-\alpha}
&=
-\frac{A}{\alpha}
\cdot
\frac{1}{\displaystyle 1-\frac{z\mathstrut}{\alpha}}
=
-\frac{A}{\alpha}
\sum_{i=0}^\infty \frac{z^i}{\alpha^i}
\end{align*}
als Potenzreihendarstellung um den Nullpunkt.
Der Konvergenzradius ist $|\alpha|$.
Eine analoge Potenzreihenentwicklung lässt sich auch für jeden beliebigen
anderen Entwicklungspunkt $z_0$ finden, solange $z_0$ keine Nullstelle
von $q(z)$ ist.

Für Nennergrad $l>1$ ist
\[
\frac{A}{(z-\alpha)^l}
=
(-1)^l
\frac{A}{\alpha^l}
\cdot
\raisebox{-6pt}{$\displaystyle
\left(
\raisebox{6pt}{$\displaystyle
\frac{1}{\displaystyle 1-\frac{z\mathstrut}{\alpha}}
$}
\right)^l
$}
=
(-1)^l
\frac{A}{\alpha^l}
\cdot
\biggl(
\sum_{i=0}^\infty \frac{z^i}{\alpha^i}
\biggr)^l.
\]
Das Produkt von Potenzreihen ist wieder eine Potenzreihe, somit
sind auch rationale Funktionen analytisch.
Ausserdem kann die Potenzreihe durch rein algebraische Umformung
gewonnen werden.
Die Tatsache, dass es sich auch um die Taylor-Reihe handelt,
spielt dabei keine Rolle.

